% Part: applied-modal-logics
% Chapter: epistemic-logic
% Section: relational-models

\documentclass[../../../include/open-logic-section]{subfiles}

\begin{document}

\olfileid{aml}{el}{rel}

\olsection{النماذج العلائقية}

المفهوم الدلالي الأساسي في منطق المعرفة هو نفسه في المنطق الجهوي المعتاد.
فلا تزال النماذج العلائقية تتألف من مجموعة عوالم وإسناد يحدد أي
ال!!{propositional variable}s تُعد «صادقة» عند أي العوالم. وإذا لم نتعامل
إلا مع فاعل واحد، كانت لدينا علاقة نفاذ واحدة كالمعتاد. أما إذا كان منطق
المعرفة متعدد الفاعلين، فتحل محل علاقة النفاذ الواحدة مجموعة من علاقات
النفاذ، واحدة لكل~$a$ في مجموعة رموز الفاعلين~$G$.

يتألف \emph{النموذج العلائقي} من مجموعة عوالم تربط بينها علاقات نفاذ
ثنائية، واحدة لكل فاعل، ومن إسناد يحدد أي ال!!{propositional variable}s
تصدق عند أي العوالم.

\begin{defn}
  \emph{نموذج} لغة المعرفة متعددة الفاعلين ثلاثي
  $\mModel{M}=\tuple{W,(R_a)_{a \in G},V}$، حيث:
  \begin{enumerate}
  \item $W$ مجموعة غير فارغة من «العوالم»؛
  \item لكل~$a \in G$، تكون~${R}_a$ علاقة نفاذ ثنائية على~$W$؛
  \item $V$ دالة تسند إلى كل~!!{propositional
    variable}~$p$ مجموعة~$V(p)$ من العوالم الممكنة.
  \end{enumerate}
  وحين يكون~$R_a ww'$، نقول إن الفاعل~$a$ \emph{يمكنه النفاذ إلى}~$w'$
  من~$w$. وحين يكون~$w \in V(p)$، نقول إن~$p$ \emph{صادق عند}~$w$.
\end{defn}

والآلية هي آلية المنطق الجهوي النظامي نفسها، مع إضافة مزيد من علاقات النفاذ
فحسب. وسنفسر في العموم علاقة النفاذ الخاصة بفاعل معطى بأنها تمثل شيئًا عن
حالته المعلوماتية. فغالبًا ما نعامل~$R_a ww'$، مثلًا، على أنه يعبر عن
اتساق~$w'$ مع معلومات~$a$ عند~$w$. أو بعبارة أخرى، لا يستطيع~$a$ عند
$w$ أن يميز بين العالمين~$w$ و$w'$.

\end{document}
